\documentclass[conference]{IEEEtran}

\usepackage[utf8]{inputenc}
\usepackage[spanish]{babel}
\usepackage{graphicx}
\usepackage{amsmath}
\usepackage{float}
\usepackage{booktabs}
\usepackage{array}
\usepackage{xcolor}
\usepackage{url}

\graphicspath{{./}}

\title{Tarea 03 -- Generaci\'on de Im\'agenes con Autoencoder:\\
PyTorch Lightning · Hydra · WandB}

\author{
\IEEEauthorblockN{Allan Bola\~nos Barrientos}
\IEEEauthorblockA{Escuela de Ingenier\'ia en Computaci\'on\\
Instituto Tecnol\'ogico de Costa Rica\\
allanbb@estudiantec.cr}
\and
\IEEEauthorblockN{Jose David Chaves Mena}
\IEEEauthorblockA{Escuela de Ingenier\'ia en Computaci\'on\\
Instituto Tecnol\'ogico de Costa Rica\\
josechaves@estudiantec.cr}
}

\begin{document}

\maketitle

\begin{abstract}
Se implementaron y compararon dos arquitecturas de autoencoder ---
VAE (\textit{Variational Autoencoder}) y U-Net AE --- para la
detecci\'on de anomal\'ias sobre el dataset MVTec AD, restringido a
cuatro clases industriales.
Cada arquitectura fue entrenada con cuatro funciones de p\'erdida
(L1, L2, SSIM, SSIM+L1), totalizando ocho experimentos controlados
mediante Hydra y monitoreados en Weights \& Biases.
Los resultados muestran que las p\'erdidas basadas en SSIM producen
mayor separaci\'on entre im\'agenes normales y defectuosas
($\Delta\mu_{\text{VAE+SSIM+L1}} = +0.00211$), mientras que L1 y L2
generan modelos que reconstruyen igualmente bien ambas clases, siendo
ineficaces para la detecci\'on de anomal\'ias.
El c\'odigo fuente y el cuaderno de experimentos est\'an disponibles
en \cite{repo} y \cite{notebook}.
\end{abstract}

% =============================================================
\section{Introducci\'on}

La detecci\'on de anomal\'ias en entornos industriales busca identificar
defectos sin ejemplos etiquetados de fallas durante el entrenamiento.
Los autoencoders son una soluci\'on natural: se entrenan \'unicamente
sobre im\'agenes normales y se espera que presenten mayor error de
reconstrucci\'on ante im\'agenes defectuosas, ya que el modelo no ha
aprendido esos patrones.

Este trabajo implementa y compara un VAE \cite{vae} y un U-Net AE sobre
el dataset MVTec AD \cite{mvtec}, usando PyTorch Lightning \cite{pytorch}
para estructurar el entrenamiento, Hydra para la gesti\'on modular de
configuraciones y WandB para el seguimiento y comparaci\'on de
experimentos.
La implementaci\'on completa se encuentra en \cite{repo} y los
experimentos reproducibles en el cuaderno \cite{notebook}.

% =============================================================
\section{Dataset}

Se utiliz\'o el dataset MVTec AD, restringido a cuatro clases:
\textbf{cable}, \textbf{capsule}, \textbf{screw} y \textbf{transistor}.
Cada imagen fue redimensionada a $128 \times 128$ p\'ixeles con 3 canales
RGB y normalizada al rango $[0,1]$.

El conjunto de entrenamiento contiene \'unicamente im\'agenes sin defectos.
El conjunto de prueba incluye muestras normales y con distintos tipos de
anomal\'ias, totalizando \textbf{542 im\'agenes}.
La partici\'on entrenamiento/validaci\'on se realiz\'o con proporci\'on
85\%/15\% sobre las muestras de entrenamiento.

% =============================================================
\section{Arquitecturas}

\subsection{VAE --- Variational Autoencoder}

El encoder convolucional tiene cinco bloques con canales
$[32,\,64,\,128,\,256,\,512]$, seguido de capas lineales que proyectan
al espacio latente gaussiano de dimensi\'on $z = 128$.
La funci\'on de p\'erdida total es:

\begin{equation}
  \mathcal{L} = \mathcal{L}_{\text{recon}}
              + \beta \cdot \mathcal{L}_{\text{KL}}, \quad \beta = 1.0
\end{equation}

\begin{equation}
  \mathcal{L}_{\text{KL}}
    = -\frac{1}{2}\sum\!\left(1 + \log\sigma^2 - \mu^2 - \sigma^2\right)
\end{equation}

\subsection{U-Net AE}

El encoder tiene cuatro bloques con canales $[64,\,128,\,256,\,512]$ y
\textit{skip connections} que conectan cada nivel del encoder con su
correspondiente en el decoder, permitiendo recuperar detalles finos de
textura.
El espacio latente es de dimensi\'on $z = 128$.
A diferencia del VAE, no impone una distribuci\'on expl\'icita sobre el
espacio latente.

% =============================================================
\section{Configuraci\'on Experimental}

\subsection{Estructura Hydra}

El proyecto se organiz\'o con la siguiente estructura de configuraci\'on:

\begin{verbatim}
conf/
  config.yaml
  model/
    vae.yaml
    u-net.yaml
  trainer/
    default.yaml
  logger/
    wandb.yaml
\end{verbatim}

Esto permite ejecutar cualquier combinaci\'on de modelo y p\'erdida
mediante un \'unico comando, por ejemplo:

\begin{verbatim}
python train.py model=vae model.loss=ssim_l1
\end{verbatim}

\subsection{Funciones de P\'erdida}

Se compararon cuatro funciones sobre cada arquitectura:

\begin{itemize}
  \item \textbf{L1}: error absoluto medio, robusto a valores at\'ipicos.
  \item \textbf{L2}: error cuadr\'atico medio, penaliza fuertemente
        errores grandes.
  \item \textbf{SSIM} \cite{ssim}: similitud estructural, modela la
        percepci\'on humana de calidad de imagen.
  \item \textbf{SSIM+L1}: combinaci\'on que aporta la estabilidad de L1
        y la sensibilidad estructural de SSIM.
\end{itemize}

\subsection{Hiperpar\'ametros Comunes}

Todos los experimentos comparten: batch size 128, learning rate
$10^{-4}$, optimizador Adam con \textit{ReduceLROnPlateau}
(factor 0.5, paciencia 8), 50 \'epocas, precisi\'on \texttt{bf16-mixed}.

% =============================================================
\section{Experimentos y Resultados}

\subsection{P\'erdida de Validaci\'on Final}

La Tabla~\ref{tab:valloss} resume la mejor p\'erdida de validaci\'on
registrada en WandB para cada configuraci\'on.
El U-Net AE converge a valores considerablemente menores en todas las
configuraciones, consecuencia directa de sus \textit{skip connections}.

\begin{table}[H]
\caption{Mejor p\'erdida de validaci\'on por configuraci\'on (50 \'epocas)}
\label{tab:valloss}
\centering
\begin{tabular}{llc}
\toprule
\textbf{Modelo} & \textbf{P\'erdida} & \textbf{Val loss} \\
\midrule
VAE   & L1      & 0.2243 \\
VAE   & L2      & 0.0897 \\
VAE   & SSIM    & 0.7293 \\
VAE   & SSIM+L1 & 0.6643 \\
\midrule
U-Net & L1      & 0.0172 \\
U-Net & L2      & 0.0008 \\
U-Net & SSIM    & 0.0189 \\
U-Net & SSIM+L1 & 0.0169 \\
\bottomrule
\end{tabular}
\end{table}

\subsection{Error de Reconstrucci\'on en Test}

La Tabla~\ref{tab:mse} resume el MSE promedio sobre las 542 im\'agenes
de prueba.
$\Delta\mu = \mu_{\text{defecto}} - \mu_{\text{normal}}$; valores positivos
indican capacidad de separaci\'on para detecci\'on de anomal\'ias.

\begin{table}[H]
\caption{MSE de reconstrucci\'on en el conjunto de prueba}
\label{tab:mse}
\centering
\resizebox{\columnwidth}{!}{%
\begin{tabular}{llccc}
\toprule
\textbf{Modelo} & \textbf{P\'erdida} &
$\mu_{\text{normal}}$ & $\mu_{\text{defecto}}$ & $\Delta\mu$ \\
\midrule
VAE   & L1      & 0.03629 & 0.02611 & \textcolor{red}{$-$0.01018} \\
VAE   & L2      & 0.01833 & 0.01366 & \textcolor{red}{$-$0.00467} \\
VAE   & SSIM    & 0.05004 & 0.03705 & \textcolor{red}{$-$0.01299} \\
VAE   & SSIM+L1 & 0.05410 & 0.05621 & \textcolor{green!60!black}{$+$0.00211} \\
\midrule
U-Net & L1      & 0.00057 & 0.00041 & \textcolor{red}{$-$0.00016} \\
U-Net & L2      & 0.00069 & 0.00052 & \textcolor{red}{$-$0.00017} \\
U-Net & SSIM    & 0.00123 & 0.00168 & \textcolor{green!60!black}{$+$0.00045} \\
U-Net & SSIM+L1 & 0.00067 & 0.00070 & \textcolor{green!60!black}{$+$0.00003} \\
\bottomrule
\end{tabular}%
}
\end{table}

\subsection{Histogramas de Error}

La Figura~\ref{fig:hist} muestra las distribuciones de MSE para los 8
modelos evaluados. Solo las configuraciones con SSIM presentan
separaci\'on positiva ($\Delta\mu > 0$) entre clases.

\begin{figure*}[t]
  \centering
  \includegraphics[width=\textwidth]{histogramas_error.png}
  \caption{Distribuciones de error de reconstrucci\'on (MSE) para im\'agenes
           normales (verde) vs.\ defectuosas (rojo) en las 542 im\'agenes
           de prueba.}
  \label{fig:hist}
\end{figure*}

La Figura~\ref{fig:defecto} desglosa el error por tipo de defecto para el
primer modelo evaluado, ilustrando la variabilidad del error seg\'un el
tipo de anomal\'ia presente.

\begin{figure}[H]
  \centering
  \includegraphics[width=\columnwidth]{error_por_defecto.png}
  \caption{Error de reconstrucci\'on (MSE) por subclase de defecto ---
           VAE con p\'erdida L1.}
  \label{fig:defecto}
\end{figure}

\subsection{Reconstrucciones Visuales}

La Figura~\ref{fig:recon} compara las reconstrucciones de validaci\'on al
final del entrenamiento para los 8 modelos. El U-Net AE produce
reconstrucciones notablemente m\'as n\'itidas gracias a las
\textit{skip connections}.

\begin{figure*}[t]
  \centering
  \includegraphics[width=\textwidth]{reconstructions_all_models.png}
  \caption{Reconstrucciones de validaci\'on en la \'ultima \'epoca.
           Fila superior: VAE (L1, L2, SSIM, SSIM+L1).
           Fila inferior: U-Net AE (mismas p\'erdidas).}
  \label{fig:recon}
\end{figure*}

\subsection{Espacio Latente (t-SNE)}

La Figura~\ref{fig:tsne} visualiza el espacio latente sobre el conjunto
de prueba mediante t-SNE. Puntos verdes representan im\'agenes normales,
rojos im\'agenes defectuosas.

\begin{figure*}[t]
  \centering
  \includegraphics[width=\textwidth]{tsne_all_models.png}
  \caption{t-SNE del espacio latente sobre el conjunto de prueba para los
           8 modelos. El VAE tiende a formar cl\'usters m\'as compactos
           gracias a la regularizaci\'on KL.}
  \label{fig:tsne}
\end{figure*}

% =============================================================
\section{An\'alisis y Discusi\'on}

\subsection{Efecto de la Funci\'on de P\'erdida}

La Tabla~\ref{tab:mse} muestra un patr\'on claro: \textbf{L1 y L2
producen $\Delta\mu$ negativo}, es decir, los modelos reconstruyen mejor
las im\'agenes defectuosas que las normales, invirtiendo el criterio de
detecci\'on.
Esto es especialmente pronunciado en U-Net AE, donde las
\textit{skip connections} permiten recuperar detalles finos incluso de
patrones no vistos durante el entrenamiento.

En contraste, \textbf{SSIM introduce sensibilidad estructural}: el modelo
aprende a reconstruir la estructura espec\'ifica de im\'agenes normales
y falla de forma medible ante anomal\'ias.
El mejor resultado de U-Net ($\Delta\mu = +0.00045$) y VAE
($\Delta\mu = +0.00211$) se obtiene con funciones que incluyen SSIM.

\subsection{VAE vs.\ U-Net AE}

El VAE presenta errores de mayor magnitud (0.03--0.05) respecto al
U-Net (0.0005--0.001) debido a la regularizaci\'on KL que comprime la
informaci\'on en el espacio latente.
Esto resulta \textit{ventajoso} para la detecci\'on de anomal\'ias:
al no poder reconstruir perfectamente, el modelo es m\'as sensible a
patrones fuera de distribuci\'on.

El U-Net AE con \textit{skip connections} logra reconstrucciones m\'as
n\'itidas en general (Figura~\ref{fig:recon}), pero esta capacidad juega
en su contra cuando se busca discriminar anomal\'ias, pues tiende a
generalizar bien incluso sobre im\'agenes defectuosas.

\subsection{Espacio Latente}

Como se aprecia en la Figura~\ref{fig:tsne}, el an\'alisis t-SNE sobre
el conjunto de prueba permite visualizar si el modelo aprendi\'o
representaciones que distinguen im\'agenes normales de defectuosas sin
supervisi\'on expl\'icita.
El VAE produce cl\'usters m\'as compactos para las im\'agenes normales,
gracias a la regularizaci\'on KL que fuerza el espacio latente a seguir
una distribuci\'on gaussiana.
En ambas arquitecturas son visibles subgrupos por tipo de producto
industrial (cable, capsule, screw, transistor).

% =============================================================
\section{Conclusiones}

\begin{itemize}
  \item \textbf{Mejor p\'erdida para VAE}: SSIM+L1 ($\Delta\mu = +0.00211$),
        combinaci\'on que aporta estabilidad de entrenamiento y
        sensibilidad estructural favorable para la detecci\'on de anomal\'ias.

  \item \textbf{Mejor p\'erdida para U-Net AE}: SSIM ($\Delta\mu = +0.00045$),
        \'unica configuraci\'on con separaci\'on positiva.
        L1 y L2 producen resultados invertidos.

  \item \textbf{Mejor configuraci\'on general}: VAE con SSIM+L1,
        con la mayor separaci\'on relativa entre clases.

  \item \textbf{Detecci\'on de anomal\'ias}: efectiva \'unicamente con
        p\'erdidas basadas en SSIM.
        L1 y L2 minimizan el error p\'ixel a p\'ixel sin capturar la
        distribuci\'on estructural de las im\'agenes normales.

  \item \textbf{Impacto de la arquitectura}: las \textit{skip connections}
        del U-Net AE mejoran la calidad de reconstrucci\'on general pero
        reducen la especificidad necesaria para detectar anomal\'ias.
\end{itemize}

\begin{thebibliography}{00}

\bibitem{mvtec}
P.~Bergmann, M.~Fauser, D.~Sattlegger, and C.~Steger,
``MVTec AD --- A Comprehensive Real-World Dataset for Unsupervised
Anomaly Detection,''
\textit{IEEE/CVF CVPR}, 2019.

\bibitem{vae}
D.~P.~Kingma and M.~Welling,
``Auto-Encoding Variational Bayes,''
\textit{ICLR}, 2014.

\bibitem{ssim}
Z.~Wang, A.~C.~Bovik, H.~R.~Sheikh, and E.~P.~Simoncelli,
``Image Quality Assessment: From Error Visibility to Structural
Similarity,''
\textit{IEEE Trans.\ Image Process.}, vol.~13, no.~4, pp.~600--612, 2004.

\bibitem{pytorch}
A.~Paszke et al.,
``PyTorch: An Imperative Style, High-Performance Deep Learning Library,''
\textit{NeurIPS 32}, 2019.

\bibitem{repo}
A.~Bola\~nos Barrientos and J.~D.~Chaves Mena,
``mvtec-autoencoder-reconstruction,''
GitHub, 2025.
[Online]. Available: \url{https://github.com/AllanDBB/mvtec-autoencoder-reconstruction}

\bibitem{notebook}
A.~Bola\~nos Barrientos and J.~D.~Chaves Mena,
``Tarea 03 --- Cuaderno de Experimentos,''
Google Colab, 2025.
[Online]. Available: \url{https://colab.research.google.com/drive/12BJ7j7Cax0TvKm4fzNAxTz6kOrY02dHn}

\end{thebibliography}

\end{document}
